% Source: Wald GR, physical PDF pp. 372--388
%         (printed pp. 359--375).
% Coverage: Section 13.2, Spinors in Curved Spacetime; equations (13.2.1)--(13.2.58).
% Figures/tables/footnotes: Figures 13.2--13.3; footnotes 6--12; no tables.
% Status: source-reviewed and compile-clean.
% Uncertainties: none; a printed equation-number error is corrected below.

\subsection{Spinors in Curved Spacetime}\label{sec:13.2}

In the previous section, we defined the notion of spinor fields on Minkowski
spacetime. Our definition was motivated by the fact that the natural action of the
Poincar\'e group on spinor fields and spinorial tensor fields gave rise to the desired
representations up to phase of the Poincar\'e group. Thus, the
\lq\lq transformation property\rq\rq\ \eqref{eq:13.1.41} under Poincar\'e
isometries was an essential ingredient of our definition of spinor fields on
Minkowski spacetime and was used to identify real spinorial tensor fields of type
\((1,0;1,0)\) with vector fields. However, the Poincar\'e group does not act in a
natural way on a curved spacetime, so clearly this characteristic property of spinor
fields cannot be carried over in a direct manner to curved spacetime. Thus, we seek
to reformulate the notion of spinor fields so that it applies in curved spacetimes
and, of course, such that it reduces in Minkowski spacetime to the notion of spinor
fields given in the previous section.

Since a general, curved spacetime possesses no isometries or any other preferred
classes of diffeomorphisms and since even in Minkowski spacetime there is no natural
action of the full group of diffeomorphisms on spinor fields, we cannot expect to
define a \lq\lq transformation law\rq\rq\ of the type~(2.2.10) under
diffeomorphisms for spinor fields in curved spacetime. However, as in Minkowski
spacetime, we may represent an observer together with his measuring apparatus at an
event, \(x\), in a curved spacetime \((M,g_{ab})\) by an orthonormal tetrad at \(x\).
Hence, associated with two different observers \(O_1\) and \(O_2\) at \(x\) is a
Lorentz transformation which rotates the tetrad of \(O_1\) into that of \(O_2\).
Note that this Lorentz transformation acts on the tangent space, \(T_xM\), rather
than the spacetime manifold, \(M\). We shall seek to define spinors at \(x\) so that
associated with each Lorentz transformation is the spinor transformation
\begin{equation}
  \psi^A(x)\longmapsto\pm L^A{}_B\psi^B(x)
  \tag{13.2.1}\label{eq:13.2.1}
\end{equation}
such that the results of all measurements by \(O_2\) on the spinor
\(\pm L^A{}_B\psi^B\) at \(x\) are identical to the results of all measurements by
\(O_1\) on \(\psi^A\). (Again, the sign ambiguity in~[13.2.1] is resolved if a
continuous curve connecting the Lorentz transformation to the identity element is
specified.) Thus, in formulating a notion of spinors in curved spacetime we shall
replace the action~\eqref{eq:13.1.41} of the Poincar\'e group of isometries on
Minkowski spacetime by the action~\eqref{eq:13.2.1} of the Lorentz group on the
tangent space at each point.

Fiber bundles provide a precise mathematical framework for defining spinor fields in
curved spacetime. We shall proceed, therefore, by defining the general notion of a
principal fiber bundle and its associated fiber bundles. The construction of the
spinor bundle then will be described.

Let \(G\) be a Lie group, let \(B\) be a manifold, and consider a \(C^\infty\) map
\(\phi:G\times B\longrightarrow B\). We shall write \(\phi(g,p)\) as
\(\phi_g(p)\). The map \(\phi\) is said to be a \emph{left action} of \(G\) on \(B\)
if (i) for each fixed \(g\in G\), the map \(\phi_g:B\longrightarrow B\) is a
diffeomorphism and (ii) for all \(g_1,g_2\in G\), we have
\(\phi_{g_1}\circ\phi_{g_2}=\phi_{g_1g_2}\). An example of a left action of \(G\)
on the manifold \(B=G\) is provided by the left translation map considered in
section~7.2. For a general left action \(\phi\) it follows from (ii) that
\(\phi_e\circ\phi_e=\phi_e\), where \(e\) is the identity element, which (composing
with \(\phi_e^{-1}\)) implies that \(\phi_e\) is just the identity map on \(B\). This
implies further that \(\phi_{g^{-1}}=\phi_g^{-1}\) for all \(g\in G\). A left action
\(\phi\) is said to be \emph{free} if for each \(g\ne e\), \(\phi_g\) leaves no point
of \(B\) fixed, i.e., if for all \(p\in B\) and \(g\ne e\) we have
\(\phi_g(p)\ne p\). [On the other hand, since \(\phi_e\) is the identity map, we have
\(\phi_e(p)=p\) for all \(p\in B\).] Thus, for example, left translation is a free
action of \(G\) on the manifold \(G\). For each \(p\in B\), the set
\(O=\{\phi_g(p)\mid g\in G\}\) is called the \emph{orbit} of \(p\) under \(G\). It is
easily seen that the condition that two points of \(B\) lie on the same orbit of
\(G\) defines an equivalence relation between points of \(B\), so \(B\) can be
expressed as a disjoint union of orbits.

In essence, a principal fiber bundle is a manifold which locally (but not necessarily
globally) \lq\lq looks like\rq\rq\ the product, \(G\times M\), of a Lie group \(G\)
and a manifold \(M\). More precisely, a \emph{principal fiber bundle}
\((B,G,M,\phi)\) consists of a manifold \(B\) (called the \emph{bundle manifold}), a
Lie group \(G\) (called the \emph{fiber group}), a manifold \(M\) (called the
\emph{base manifold}), and a free left action
\(\phi:G\times B\longrightarrow B\) satisfying the following two properties: (i)
The orbits of \(G\) are in one-to-one, onto correspondence with the points of
\(M\), and the projection map \(\pi:B\longrightarrow M\) which assigns to each
\(p\in B\) the point of \(M\) associated with the orbit of \(p\) is \(C^\infty\).
(ii) For each \(x\in M\) there exists an open neighborhood, \(U\), of \(x\) such
that there is a diffeomorphism, \(\psi\), taking \(\pi^{-1}(U)\subseteq B\) into
\(G\times U\) such that the action of \(G\) on \(\pi^{-1}(U)\) corresponds to left
multiplication on \(G\times U\); i.e., if \(\psi(p)=(g,x)\), then
\(\psi[\phi_{g'}(p)]=(g'g,x)\). Figure~13.2 illustrates the nature of a principal
fiber bundle. Note that we always have
\(\dim(B)=\dim(G)+\dim(M)\).

% Source: Wald GR, physical PDF p. 374 (printed p. 361).
% Coverage: Figure 13.2, principal fiber bundle.
% Figures/tables/footnotes: one figure asset; no tables or footnotes.
% Status: source-reviewed and compile-clean.
% Uncertainties: none.

\begingroup
\renewcommand{\thefigure}{13.2}
\begin{figure}[htbp]
  \centering
  \begin{tikzpicture}[scale=.85,>=Stealth]
    \draw[thick] (-3.6,1.15) rectangle (3.6,4.15);
    \node at (2.55,3.75) {\(B\)};
    \draw[dashed] (-1.2,1.15) -- (-1.2,4.15);
    \draw[dashed] (.7,1.15) -- (.7,4.15);
    \draw[thick] (-.15,1.15) .. controls (.17,2.05) and (-.32,2.9) .. (.08,4.15);
    \draw[->,thick] (1.30,1.15) -- (1.30,-.64) node[midway,right] {\(\pi\)};
    \draw[thick] (-3.1,-.7) -- (3.1,-.7);
    \node[below] at (-3.1,-.7) {\(M\)};
    \fill (0,-.7) circle (1.5pt);
    \node[below] at (0,-.7) {\(x\)};
    \draw[dashed] (-.75,-.45) -- (-.75,-.95);
    \draw[dashed] (.75,-.45) -- (.75,-.95);
    \draw[->] (.9,-1.2) .. controls (.68,-1.02) and (.38,-1.0) .. (.18,-.82);
    \node[right] at (.92,-1.2) {\(U\)};
    \draw[->] (-2.0,2.55) .. controls (-1.42,2.5) and (-.78,2.27) .. (-.28,2.05);
    \node[left] at (-2.05,2.6) {\(\pi^{-1}[x]\)};
    \draw[->] (2.3,2.08) .. controls (1.62,2.1) and (1.12,2.34) .. (.87,2.55);
    \node[right] at (2.35,2.1) {\(\pi^{-1}[U]\)};
    \draw[thick] (-5.2,1.2) -- (-5.2,4.1);
    \node[left] at (-5.2,2.7) {\(G\)};
  \end{tikzpicture}
  \caption{A diagram illustrating the structure of a principal fiber bundle (see
  text).}
  \label{fig:13.2}
\end{figure}
\endgroup


Thus, a particularly simple example of a principal fiber bundle is obtained by
taking the product manifold, \(B=G\times M\), of a Lie group \(G\) and a manifold
\(M\), with the left action of \(G\) on \(B\) defined by left multiplication, i.e.,
\(\phi_{g'}(g,x)=(g'g,x)\). A bundle of this form is said to be \emph{trivial}. One
of the simplest examples of a nontrivial principal fiber bundle is obtained by taking
\(B\) to be the circle, \(S^1\), and \(G\) to be the group, \(Z_2\), consisting of
the two elements \(e,a\) with \(a^2=e\). (Thus, \(G\) is a zero-dimensional,
disconnected Lie group.) We define a left action of \(Z_2\) on the circle,
\(\phi:Z_2\times B\longrightarrow B\), by \(\phi_e(\theta)=\theta\) and
\(\phi_a(\theta)=\theta+\pi\). Thus, the orbits of \(Z_2\) are the opposite points
on the circle \(B\), and the collection of orbits, \(M\), can be given the manifold
structure \(S^1\) so that the projection \(\pi:B\longrightarrow M\) is smooth. It
then may be verified that \((S^1,Z_2,S^1,\phi)\) is a principal fiber bundle. This
bundle is nontrivial since \(B\) is not diffeomorphic to the product manifold
\(G\times M\), which consists of two disconnected circles. The bundles \(B\) and
\(G\times M\) are illustrated in Figure~13.3.

% Source: Wald GR, physical PDF p. 375 (printed p. 362).
% Coverage: Figure 13.3, trivial and nontrivial \(Z_2\) principal bundles.
% Figures/tables/footnotes: one figure asset; no tables or footnotes.
% Status: source-reviewed and compile-clean.
% Uncertainties: none.

\begingroup
\renewcommand{\thefigure}{13.3}
\begin{figure}[htbp]
  \centering
  \begin{tikzpicture}[scale=.8,>=Stealth]
    \begin{scope}[shift={(-4,0)}]
      \node at (-2.65,2.25) {\(G=Z_2\)};
      \fill (-1.55,2.50) circle (2pt);
      \fill (-1.55,2.00) circle (2pt);
      \draw[->] (-2.05,2.28) -- (-1.63,2.47);
      \draw[->] (-2.05,2.28) -- (-1.63,2.04);
      \draw[thick] (0,2.55) ellipse (1.05 and .35);
      \draw[thick] (0,1.65) ellipse (1.05 and .35);
      \fill (0,2.20) circle (1.7pt);
      \fill (0,1.30) circle (1.7pt);
      \draw[->,thick] (0,1.25) -- (0,.45) node[midway,right] {\(\pi\)};
      \draw[thick] (0,0) ellipse (1.05 and .35);
      \fill (0,-.35) circle (1.7pt);
      \node[right] at (.7,-.38) {\(M=S^1\)};
      \node[below] at (0,-1.05) {\((a)\)};
    \end{scope}
    \begin{scope}[shift={(4,0)}]
      \node at (-2.65,2.25) {\(G=Z_2\)};
      \fill (-1.55,2.50) circle (2pt);
      \fill (-1.55,2.00) circle (2pt);
      \draw[->] (-2.05,2.28) -- (-1.63,2.47);
      \draw[->] (-2.05,2.28) -- (-1.63,2.04);
      % One circle, drawn in projection with a self-crossing; two intersecting
      % ellipses would incorrectly depict a disconnected pair of fibers.
      \draw[thick,smooth,variable=\t,domain=0:360,samples=80]
        plot ({1.15*sin(\t)},{2.10+.43*sin(2*\t)});
      % Break the under-strand at the projected crossing so the drawing reads
      % as one immersed circle, rather than as two circles joined at a vertex.
      \draw[thick,white,line width=3.2pt,smooth,variable=\t,domain=171:189,samples=12]
        plot ({1.15*sin(\t)},{2.10+.43*sin(2*\t)});
      \draw[thick,smooth,variable=\t,domain=171:189,samples=12]
        plot ({1.15*sin(\t)},{2.10+.43*sin(2*\t)});
      % A Z_2 orbit consists of two antipodal points of this single circle.
      \fill (-.996,2.472) circle (1.7pt);
      \fill ( .996,2.472) circle (1.7pt);
      \draw[->,thick] (0,1.35) -- (0,.45) node[midway,right] {\(\pi\)};
      \draw[thick] (0,0) ellipse (1.05 and .35);
      \fill (0,-.35) circle (1.7pt);
      \node[right] at (.7,-.38) {\(M=S^1\)};
      \node[below] at (0,-1.05) {\((b)\)};
    \end{scope}
  \end{tikzpicture}
  \caption{(a) The trivial principal fiber bundle \(Z_2\times S^1\). (b) The
  principal bundle \((S^1,Z_2,S^1,\phi)\) constructed in the text.}
  \label{fig:13.3}
\end{figure}
\endgroup


A particularly important example of a principal fiber bundle is the \emph{bundle of
bases}, defined as follows. Let \(M\) be an \(n\)-dimensional manifold and consider
the collection, \(B\), of pairs \((x,(e_\mu)^a)\) where \(x\in M\) and
\((e_\mu)^a\) (where \(\mu=1,\ldots,n\)) is a basis of the tangent space,
\(T_xM\). We can put a manifold structure on \(B\) as follows. If
\(\psi:U\longrightarrow\mathbb R^n\) is a chart on a region, \(U\), of \(M\), we
obtain a coordinate basis \((\partial/\partial x^\mu)^a\) associated with \(\psi\).
At each \(x\in U\) we can expand each \((e_\mu)^a\) in terms of this basis,
\((e_\mu)^a=\sum_\nu\alpha_\mu{}^\nu(\partial/\partial x^\nu)^a\). Let
\(\widetilde U\) consist of all points \((x,(e_\mu)^a)\) of \(B\) with \(x\in U\).
We define \(\widetilde\psi:\widetilde U\longrightarrow\mathbb R^{n+n^2}\) by
\begin{equation}
  \widetilde\psi(x,(e_\mu)^a)=(x^\sigma,\alpha_\mu{}^\nu).
  \tag{13.2.2}\label{eq:13.2.2}
\end{equation}
We use the collection of all maps of this form to define a family of charts on
\(B\). It is easily checked that these charts are compatible, thus making \(B\) a
manifold of dimension \(n+n^2\). The general linear group in \(n\) dimensions,
\(\mathrm{GL}(n)\)---which consists of the invertible linear maps taking
\(\mathbb R^n\) into itself with group multiplication defined by
composition---acts on the left on \(B\) as follows. We view each
\(A\in\mathrm{GL}(n)\) as a matrix, \(A^\mu{}_\nu\), of components of the linear map
\(A\) in the natural basis of \(\mathbb R^n\). We define
\(\phi_A:B\longrightarrow B\) by
\begin{equation}
  \phi_A(x,(e_\mu)^a)
  =\left(x,\sum_\nu(A^{-1})^\nu{}_\mu(e_\nu)^a\right);
  \tag{13.2.3}\label{eq:13.2.3}
\end{equation}
i.e., we use the inverse map, \(A^{-1}\), to transform the basis of the tangent
space at \(x\). Then it is easily seen that \(\phi\) is a free left action of
\(\mathrm{GL}(n)\) on \(B\) and that the above properties (i) and (ii) are
satisfied. Thus, \((B,\mathrm{GL}(n),M,\phi)\) is a principal fiber bundle. For
some manifolds, \(M\)---for example, for \(\mathbb R^n\)---the bundle manifold has
the structure \(\mathrm{GL}(n)\times M\), i.e., the bundle of bases is trivial. On
the other hand, the bundle of bases of other manifolds, such as the 2-sphere,
\(S^2\), is nontrivial.

Similarly, for a manifold \(M\) on which a metric, \(g_{ab}\), of signature
\((p,q)\) is defined, we may construct the \emph{bundle of orthonormal bases}. The
only changes from the above construction are that we take \(B\) to consist of pairs
\((x,(e_\mu)^a)\), where \(x\in M\) and \((e_\mu)^a\) now is an orthonormal basis for
\(T_xM\), and the group which acts on \(B\) is now the orthogonal group \(O(p,q)\),
rather than \(\mathrm{GL}(n)\). In particular, for a spacetime \((M,g_{ab})\) the
bundle of orthonormal bases has fiber group equal to the (improper) Lorentz group
\(O(3,1)\). Similarly, for an orientable, time orientable spacetime, we may
construct the bundle of oriented, time oriented bases with the proper Lorentz group,
\(\Lambda\), as the fiber group.

Given a principal fiber bundle \((B,G,M,\phi)\), another manifold \(F\), and a (not
necessarily free) left action \(\chi:G\times F\longrightarrow F\) of \(G\) on
\(F\), there is a general procedure for constructing from \(B\) a new manifold
\(E\), in which, in effect, the fiber group \(G\) is replaced by the manifold \(F\).
To do so, we take the product manifold \(B\times F\) and define a left action,
\(\Psi\), of \(G\) on \(B\times F\) by
\begin{equation}
  \Psi_g[b,f]=[\phi_g(b),\chi_g(f)].
  \tag{13.2.4}\label{eq:13.2.4}
\end{equation}
We define \(E\) to be the set of orbits of \(G\) on \(B\times F\). There is a
natural projection map \(p:E\longrightarrow M\) given by \(p(y)=\pi(b)\), where
\(b\in B\) is such that \((b,f)\in B\times F\) lies on the orbit of \(y\in E\) and
\(\pi:B\longrightarrow M\) is the projection map on the principal fiber bundle. For
every neighborhood \(U\subseteq M\) satisfying property~(ii) of the definition of
principal fiber bundle, it follows that \(p^{-1}(U)\subseteq E\) is homeomorphic to
\(F\times U\). We define a manifold structure on \(E\) by requiring that these
homeomorphisms be diffeomorphisms. Thus, \(E\) locally \lq\lq looks like\rq\rq\
\(F\times M\) in the same sense as \(B\) locally \lq\lq looks like\rq\rq\
\(G\times M\). The manifold \(E\) together with the large amount of structure on
\(E\) determined by \(B,G,M,\phi,F\), and \(\chi\) is called a \emph{fiber bundle}
or, more precisely, the fiber bundle associated to \((B,G,M,\phi)\) with fiber
manifold \(F\) and group action \(\chi\). For each \(x\in M\), the subset
\(p^{-1}(x)\subseteq E\)---called the \emph{fiber over \(x\)}---is diffeomorphic to
\(F\).

It may appear from the above rather complicated constructions and definitions that
fiber bundles comprise an extremely specialized class of manifolds. In fact,
however, a large variety of manifolds can be expressed in a natural and very useful
way as fiber bundles. The utility of the fiber bundle viewpoint for proving theorems
on the topology of manifolds can be seen, for example, in Steenrod (1951).

Given a principal fiber bundle \((B,G,M,\phi)\) we may take \(F=G\) and \(\chi_g\)
to be left translation. The resulting associated fiber bundle \(E\) is
diffeomorphic to \(B\), so every principal fiber bundle also can be viewed as an
associated fiber bundle. Another particularly simple example of a fiber bundle
associated to a principal bundle \((B,G,M,\phi)\) is obtained by letting \(F\) be
any manifold and taking \(\chi\) to be the trivial action,
\(\chi_g(f)=f\) for all \(g\in G\) and \(f\in F\). The resulting fiber bundle \(E\)
is diffeomorphic to the product manifold, \(F\times M\), of the fiber manifold
\(F\) with the base manifold \(M\). A fiber bundle which is diffeomorphic to
\(F\times M\) is called \emph{trivial}. A simple example of a nontrivial fiber
bundle, associated to the principal bundle \((S^1,Z_2,S^1,\phi)\) discussed above,
with fiber manifold \(F=\mathbb R\) is obtained from the following action of
\(Z_2=\{e,a\}\) on \(\mathbb R\). We define \(\chi_e(x)=x\),
\(\chi_a(x)=-x\) for all \(x\in\mathbb R\). The resulting fiber bundle \(E\)
\lq\lq locally looks like\rq\rq\ the cylinder \(\mathbb R\times S^1\), but has the
property that the \(\mathbb R\)-fibers \lq\lq flip upside down\rq\rq\ when one goes
once around a curve which projects down to a circle in the base manifold \(M=S^1\).
The manifold \(E\) is called the \emph{M\"obius strip} and also could be constructed
by identifying the points in the plane, \(\mathbb R^2\), via
\((x,y)=(x+1,-y)\).

An important example of a fiber bundle associated to the bundle of bases of an
\(n\)-dimensional manifold \(M\) is obtained as follows. We take \(F=\mathbb R^n\)
and let \(\mathrm{GL}(n)\) act on \(\mathbb R^n\) in the natural way, i.e., for
\(a=(a^1,\ldots,a^n)\in\mathbb R^n\) and \(A\in\mathrm{GL}(n)\) we define
\begin{equation}
  \bigl(\chi_A(a)\bigr)^\mu=\sum_{\nu=1}^nA^\mu{}_\nu a^\nu.
  \tag{13.2.5}\label{eq:13.2.5}
\end{equation}
Now, corresponding to each point \((x,(e_\mu)^a,a^\nu)\) in the manifold
\(B\times\mathbb R^n\) we can associate the vector
\[
  v^a=\sum_{\mu=1}^n a^\mu(e_\mu)^a
\]
at \(x\in M\). Given the actions~\eqref{eq:13.2.3} and~\eqref{eq:13.2.5} of \(G\)
on \(B\) and on \(\mathbb R^n\), it is easily seen that two points of
\(B\times\mathbb R^n\) correspond to the same vector \(v^a\) at \(x\) if and only if
they lie on the same orbit~\eqref{eq:13.2.4}. Thus, the points of the associated
fiber bundle \(E\) are in one-to-one correspondence with the pairs \((x,v^a)\),
where \(x\in M\) and \(v^a\in T_xM\). We call \(E\) the \emph{tangent bundle} of
\(M\) and denote it by \(T(M)\). If \(T(M)\) is trivial, we say that \(M\) is
\emph{parallelizable}. The bundle \(T_{k,l}(M)\) of tensors of type \((k,l)\) over
\(M\) can be constructed similarly by taking \(F\) to be the tensor product,
\(F_{k,l}\), of \(k\) copies of \(\mathbb R^n\) and \(l\) copies of its dual space
\((\mathbb R^n)^\vee\) and taking \(\chi\) to be the natural action of
\(\mathrm{GL}(n)\) on \(F_{k,l}\) given by
\begin{align}
  \bigl(\chi_A(b)\bigr)^{\mu_1\cdots\mu_k}{}_{\nu_1\cdots\nu_l}
  ={}&\sum_{\alpha_i,\beta_j=1}^n
  A^{\mu_1}{}_{\alpha_1}\cdots A^{\mu_k}{}_{\alpha_k}
  (A^{-1})^{\beta_1}{}_{\nu_1}\cdots
  (A^{-1})^{\beta_l}{}_{\nu_l}\,
  b^{\alpha_1\cdots\alpha_k}{}_{\beta_1\cdots\beta_l}.
  \tag{13.2.6}\label{eq:13.2.6}
\end{align}
If a metric, \(g_{ab}\), of signature \((p,q)\) is given on \(M\) we also may
construct \(T_{k,l}(M)\) by restricting attention to orthonormal bases, i.e., by
starting with the principal bundle of orthonormal bases and defining an action of
\(O(p,q)\) on \(F_{k,l}\) via equation~\eqref{eq:13.2.6}.

Thus, the above construction gives us a new viewpoint on tangent vectors: A tangent
vector at a point \(x\) on a manifold \(M\) is a point of the tangent bundle
\(T(M)\) lying in the fiber over \(x\), i.e., a point \(y\in T(M)\) such that
\(p(y)=x\). More generally, a tensor of type \((k,l)\) at \(x\in M\) is a point in
the fiber over \(x\) of the bundle \(T_{k,l}(M)\) of tensors of type \((k,l)\). A
smooth tensor field on \(M\) is a \emph{cross section} of \(T_{k,l}(M)\), i.e., a
\(C^\infty\) map \(\Sigma:M\longrightarrow T_{k,l}(M)\) such that
\(p\circ\Sigma\) is the identity map on \(M\). Note that this viewpoint on tensors
is rather close in spirit to the characterization of tensors mentioned in chapter~2
in terms of the tensor transformation law~(2.3.8). There it was remarked that the
change in the coordinate basis components of a tensor under a coordinate
transformation could be used to characterize it. Here, a point in \(B\times F_{k,l}\)
is a collection of \(n^{k+l}\) numbers associated with a point \(x\in M\) and a (not
necessarily coordinate) basis of the tangent space of \(x\). The group action
\eqref{eq:13.2.6} can be viewed as telling us how this collection of numbers changes
when we make a change of basis. The equivalence class of the collection of numbers
under changes of basis at \(x\) is a point of \(T_{k,l}(M)\), i.e., a tensor at \(x\).
Both the tensor transformation law and the fiber bundle characterization of tensors
have little advantage over the direct definition of tensors as multilinear maps given
in chapter~2. However, since there is no analogous direct definition of spinors, we
must proceed to define spinors by the relatively indirect route of specifying their
behavior under \(\mathrm{SL}(2,\mathbb C)\) transformations associated with changes
of orthonormal basis. We turn, now, to this task.

The basic idea for constructing spinor fields in a curved spacetime \((M,g_{ab})\)
is to start with the principal fiber bundle of oriented, time oriented orthonormal
bases, so that each fiber is diffeomorphic to the proper Lorentz group. Then we
\lq\lq unwrap\rq\rq\ each fiber to produce a principal
\(\mathrm{SL}(2,\mathbb C)\) bundle over \(M\). The spinor bundle then is
constructed as the fiber bundle associated to this principal bundle with fiber
\(F=\mathbb C^2\), where \(\chi\) is taken to be the natural action of
\(\mathrm{SL}(2,\mathbb C)\) on \(\mathbb C^2\). A spinor at \(x\in M\) may then be
defined as a point in the fiber over \(x\) of the spinor bundle.

However, complications may arise in this construction if the topology of the
spacetime \(M\) is nontrivial. Consider, first, the case where \(M\) is simply
connected. Then \(M\) is orientable and time orientable, so there is no difficulty
in constructing the principal bundle, \((B,\Lambda,M,\phi)\) of oriented, time
oriented bases. Consider, now, a closed curve, \(\gamma\), in \(B\). Since \(M\) is
simply connected, the curve \(\pi\circ\gamma\) obtained by projecting \(\gamma\)
down onto \(M\) can be continuously deformed to the trivial curve through \(x\),
i.e., the curve \(e(t)=x\) for all \(t\). This implies that in \(B\) the curve
\(\gamma\) can be continuously deformed to a curve which lies entirely in the fiber
\(\pi^{-1}[x]\) over \(x\). However, this fiber is diffeomorphic to the proper
Lorentz group, \(\Lambda\), and \(\pi_1(\Lambda)=Z_2\). Thus, in general, we would
expect there to exist precisely two distinct homotopy equivalence classes of closed
curves in \(B\). However, it is possible for there to exist only one homotopy
equivalence class of curves, i.e., for \(B\) to be simply connected. Although the
curve in the fiber over \(x\) corresponding to a \(2\pi\) rotation of the basis of
\(T_xM\) cannot be deformed to the trivial curve while remaining within the fiber over
\(x\), for certain manifolds this curve \emph{can} be continuously deformed to the
trivial curve in \(B\) by moving it out of the fiber over \(x\). In other words, it
may happen that a \(2\pi\)-rotation of a tetrad can be \lq\lq undone\rq\rq\ by
transporting the tetrad along a one-parameter sequence of closed curves through
\(x\) in \(M\). \emph{If \(B\) is simply connected, the notion of spinors on \(M\)
cannot be defined.}\footnote{However, one still may define \lq\lq generalized spin
structures\rq\rq; see Avis and Isham (1980) and the references cited therein.} This
is because we cannot consistently assign a change of sign to a spinor under a
\(2\pi\)-rotation of the basis if this \(2\pi\)-rotation can be continuously deformed
to the identity. Examples of manifolds for which \(B\) is simply connected are given
by Geroch (1968a). It is known (Milnor 1963; Clarke 1971) that \(B\) will fail to be
simply connected---i.e., \(\pi_1(B)=Z_2\) and spinors can be defined---if and only
if the second Stiefel--Whitney class (defined, e.g., in Steenrod 1951) of \(M\)
vanishes. Geroch (1968a) has proven\footnote{Since \(O(n)\) and \(O(n,1)\) have
fundamental group \(Z_2\) for all \(n>2\), one can define analogous notions of
spinors for Riemannian and Lorentzian spaces of dimension greater, respectively,
than 2 and 3. (See problem~1 for the Riemannian case with \(n=3\); see Yip (1983)
for a discussion of spinors in two-dimensional spacetimes, where \(O(1,1)\) is
simply connected.) Geroch's parallelizability criterion applies only to
four-dimensional spacetimes.} that if \(M\) is noncompact, then
\(\pi_1(B)=Z_2\) if and only if \(M\) is parallelizable. Further equivalent criteria
are given by Geroch (1968a, 1970a) and Clarke (1971).

If \(\pi_1(B)=Z_2\), we define spinors on \(M\) as follows. The universal covering
manifold, \(\widehat B\), of \(B\) will have the natural structure of a principal
fiber bundle with fiber group \(\mathrm{SL}(2,\mathbb C)\). We define the
\emph{spinor bundle}, \(S(M)\), to be the fiber bundle associated to
\((\widehat B,\mathrm{SL}(2,\mathbb C),M,\widehat\phi)\) with fiber manifold
\(\mathbb C^2\), with the action \(\chi_L\) of
\(L\in\mathrm{SL}(2,\mathbb C)\) on \(c\in\mathbb C^2\) taken to be the natural
action given by
\begin{equation}
  [\chi_L(c)]^\Sigma=\sum_{\Gamma=1}^2L^\Sigma{}_\Gamma c^\Gamma.
  \tag{13.2.7}\label{eq:13.2.7}
\end{equation}
A \emph{spinor} at \(x\in M\) is defined to be a point of \(S(M)\) lying in the
fiber over \(x\). Similarly, the bundle,
\(S_{k,l;k',l'}(M)\) of spinorial tensors of type \((k,l;k',l')\) is defined by
taking \(F\) to be the tensor product of \(k\) copies of \(\mathbb C^2\), \(l\)
copies of its dual space \((\mathbb C^2)^\vee\), \(k'\) copies of its complex conjugate
space \(\overline{\mathbb C}^{\,2}\), and \(l'\) copies of its complex conjugate
dual space \((\overline{\mathbb C}^{\,2})^\vee\). The action of \(\chi\) on
\[
  c^{\Sigma_1\cdots\Sigma_k\Sigma_1'\cdots\Sigma_{k'}'}
    {}_{\Lambda_1\cdots\Lambda_l\Lambda_1'\cdots\Lambda_{l'}'}\in F
\]
is given by
\begin{align}
  [\chi_L(c)]^{\Sigma_1\cdots\Sigma_k\Sigma_1'\cdots\Sigma_{k'}'}
    {}_{\Lambda_1\cdots\Lambda_l\Lambda_1'\cdots\Lambda_{l'}'}
  ={}&\sum_{\substack{
          \Gamma_1\cdots\Gamma_k,\Delta_1\cdots\Delta_l,\\[-.25ex]
          \Gamma_1'\cdots\Gamma_{k'}',\Delta_1'\cdots\Delta_{l'}'}}\notag\\[-.5ex]
  &{}\quad\times L^{\Sigma_1}{}_{\Gamma_1}\cdots
    L^{\Sigma_k}{}_{\Gamma_k}
    (L^{-1})^{\Delta_1}{}_{\Lambda_1}\cdots\notag\\
  &{}\quad\times (L^{-1})^{\Delta_l}{}_{\Lambda_l}
    \overline L^{\Sigma_1'}{}_{\Gamma_1'}\cdots
    \overline L^{\Sigma_{k'}'}{}_{\Gamma_{k'}'}\notag\\
  &{}\quad\times(\overline L^{-1})^{\Delta_1'}{}_{\Lambda_1'}\cdots
    (\overline L^{-1})^{\Delta_{l'}'}{}_{\Lambda_{l'}'}\notag\\
  &{}\quad\times
  c^{\Gamma_1\cdots\Gamma_k\Gamma_1'\cdots\Gamma_{k'}'}
    {}_{\Delta_1\cdots\Delta_l\Delta_1'\cdots\Delta_{l'}'}.
  \tag{13.2.8}\label{eq:13.2.8}
\end{align}
Equivalently, once the notion of spinors has been given, spinorial tensors may be
defined in terms of linear (and antilinear) maps on spinors in a manner analogous to
the way ordinary tensors were constructed from vectors in chapter~2. As in the
previous section we shall use capital Latin indices to denote spinorial tensors.

The above definition of spinors is sufficiently abstruse that some words of
explanation might be helpful. Recall, first, that a point of \(B\) consists of a
point \(x\in M\) together with an oriented, time oriented orthonormal basis at \(x\).
If we fix a particular basis \(\{(e_\mu)^a\}\) at \(x\), a point of \(\widehat B\) may
be viewed as consisting of a point \(x\in M\) together with a continuous,
one-parameter family of oriented, time oriented bases at \(x\) starting from
\(\{(e_\mu)^a\}\), where two such families are considered equivalent if they have the
same endpoint and are homotopic. In other words, the fiber of \(\widehat B\) over
\(x\) can be viewed as consisting of all oriented, time oriented orthonormal bases
at \(x\) together with their \(2\pi\)-rotations. A point of
\(\widehat B\times\mathbb C^2\) is a pair of complex numbers associated with \(x\)
and a basis or \(2\pi\)-rotated basis at \(x\). The group action \(\chi_g\) of
equation~\eqref{eq:13.2.7} tells us how these numbers
\lq\lq transform\rq\rq\ under a change of basis. Each equivalence class of
transformed numbers and bases at \(x\)---i.e., each orbit of \(\Psi_g\),
equation~\eqref{eq:13.2.4}---defines a spinor at \(x\). Thus, our fiber bundle
construction is just a precise way of implementing the notion that \lq\lq a spinor at
\(x\) is an ordered pair of complex numbers which transforms by the natural
representation of \(\mathrm{SL}(2,\mathbb C)\) under a change of basis.\rq\rq\

Note that \(\mathbb C^2\) has a great deal of structure beyond that of a complex
two-dimensional vector space; for example, it has a natural inner product and a
notion of the real and imaginary parts of vectors. However, only those properties of
\(\mathbb C^2\) which are preserved under the action of
\(\mathrm{SL}(2,\mathbb C)\) survive to yield structure on the fibers of the spinor
bundle \(S(M)\). Thus, for example, since there are elements of
\(\mathrm{SL}(2,\mathbb C)\) which take real vectors in \(\mathbb C^2\) to complex
vectors, there is no natural notion of the real and imaginary parts of a spinor at
\(x\). However, since the \(\mathrm{SL}(2,\mathbb C)\) maps are linear, they
preserve addition and scalar multiplication, so the spinors at \(x\) do have the
natural structure of a two-dimensional complex vector space. Beyond that, only the
element \(\epsilon_{\Sigma\Gamma}\in(\mathbb C^2)^\vee\otimes(\mathbb C^2)^\vee\) (as well
as objects constructed from it) given by
\[
  \epsilon_{\Sigma\Gamma}
  =\begin{pmatrix}0&1\\-1&0\end{pmatrix}
\]
is preserved under the action of \(\mathrm{SL}(2,\mathbb C)\), so the only
additional natural structure of spinors is that of a tensor \(\epsilon_{AB}\) of
type \((0,2;0,0)\). Thus, each fiber of \(S(M)\) has precisely the natural
structure of spinor space \((W,\epsilon_{AB})\). Note that there is no natural way
of identifying the different fibers of \(S(M)\), i.e., there is no way of saying
that a spinor at \(x\in M\) is \lq\lq the same\rq\rq\ as one at \(y\in M\), just as
there is no way of saying that tangent vectors at different points are
\lq\lq the same.\rq\rq\ However, we do know what it means for a spinor to vary
smoothly from point to point: A smooth spinor field is simply a (smooth) cross
section of \(S(M)\).

The relation between spinorial tensor fields of type \((1,0;1,0)\) and vector
fields may be seen as follows. Fix an isomorphism, \(\sigma\), between the real
elements of \(\mathbb C^2\otimes\overline{\mathbb C}^{\,2}\) and \(\mathbb R^4\)---such
as that given by equation~\eqref{eq:13.1.30}---for which the Lorentz metric
\(\epsilon_{\Sigma\Lambda}\overline\epsilon_{\Sigma'\Lambda'}\) on
\(\operatorname{Re}(\mathbb C^2\otimes\overline{\mathbb C}^{\,2})\) is taken into
the Lorentz metric \(\operatorname{diag}(1,-1,-1,-1)\) on \(\mathbb R^4\). Then the
natural 2 to 1 map of \(\widehat B\) onto \(B\) gives rise to a 2 to 1 map
\(\sigma'\) of
\(\widehat B\times\operatorname{Re}(\mathbb C^2\otimes\overline{\mathbb C}^{\,2})\)
onto \(B\times\mathbb R^4\). Furthermore, \(\sigma'\) commutes with the natural
group actions \(\Psi_L\) and \(\Psi_\Lambda\) on these spaces, i.e., we have
\(\sigma'\circ\Psi_L=\Psi_\Lambda\circ\sigma\), where \(\Lambda\) is the Lorentz
group element associated with \(L\in\mathrm{SL}(2,\mathbb C)\). Hence, we obtain a
(one-to-one) correspondence, \(\widetilde\sigma\), between the
\(\mathrm{SL}(2,\mathbb C)\) orbits of
\(\widehat B\times\operatorname{Re}(\mathbb C^2\otimes\overline{\mathbb C}^{\,2})\)
and the Lorentz group orbits of \(B\times\mathbb R^4\). Thus, we obtain in turn an
identification of the real subset of \(S_{1,0;1,0}(M)\) with \(T(M)\); i.e., real
spinorial tensors of type \((1,0;1,0)\) are identified with vectors. As in the
previous section, we shall incorporate this identification into our notation by
allowing \(v^{AA'}\) to denote a vector rather than explicitly writing
\(\widetilde\sigma^{\,a}{}_{AA'}v^{AA'}\). Again the spacetime metric is related to
\(\epsilon_{AB}\) by
\begin{equation}
  g_{AA'BB'}=\epsilon_{AB}\overline\epsilon_{A'B'}.
  \tag{13.2.9}\label{eq:13.2.9}
\end{equation}

In the case where the spacetime manifold, \(M\), is not simply connected, the
analysis of whether the notion of spinors can be defined is somewhat more
complicated. First, \(M\) may fail to be orientable or time oriented. \emph{In such
a case, the bundle, \(B\), of oriented, time oriented orthonormal bases does not
exist, and the notion of spinors cannot be defined.}\footnote{We should reemphasize
that the term \lq\lq spinors\rq\rq\ here and elsewhere in this chapter means
\(\mathrm{SL}(2,\mathbb C)\) spinors. Other types of spinors can be defined by a
similar construction starting from other bundles of bases. In particular, in a time
orientable but not necessarily orientable spacetime, let \(\widetilde B\) denote the
bundle of time oriented orthonormal bases. The fiber group is then the
(disconnected) group obtained by composing proper Lorentz transformations with a
\lq\lq parity\rq\rq\ transformation. Its covering group can be made to act on a
four-dimensional complex vector space via the usual transformation formulas for
Dirac spinors. Thus, a notion of (4-component) Dirac spinors can be defined on
nonorientable spacetimes. In the case of an orientable spacetime, Dirac spinors are
naturally isomorphic to a pair \((\phi^A,\sigma_{A'})\) consisting of a
two-component spinor \(\phi^A\) and a complex conjugate dual spinor
\(\sigma_{A'}\), as mentioned at the end of section~13.1.} If \(M\) is orientable
and time orientable, we may construct \(B\), but we cannot obtain \(\widehat B\) by
taking the universal covering manifold of \(B\) since this would \lq\lq unwrap\rq\rq\
\(M\) as well as the fibers of \(B\). In order to make sense of the idea of
unwrapping only the fibers of \(B\), it is necessary that the fundamental group of
\(B\) be of the direct product form
\begin{equation}
  \pi_1(B)=G_1\times G_2,
  \tag{13.2.10}\label{eq:13.2.10}
\end{equation}
where \(G_1\) is the fundamental group, \(\pi_1(\Lambda)=Z_2\), of the Lorentz
group and \(G_2\) is the fundamental group, \(\pi_1(M)\), of \(M\). [Here the direct
product \(G=G_1\times G_2\) of groups \(G_1\) and \(G_2\) is defined as the group
consisting of ordered pairs \((g_1,g_2)\) with \(g_1\in G_1\), \(g_2\in G_2\) with
composition law \((g_1,g_2)(g_1',g_2')=(g_1g_1',g_2g_2')\).] Furthermore, it is
understood in equation~\eqref{eq:13.2.10} that each element of the subgroup of
\(\pi_1(B)\) of the form \((g_1,e_2)\) with \(e_2\) the identity element of \(G_2\)
is homotopic to the closed curve \(g_1\) lying within a single Lorentz group fiber,
whereas each element of the subgroup \((e_1,g_2)\) projects down to the closed curve
\(g_2\in\pi_1(M)\) in \(M\). Note that in the case where \(M\) is simply connected,
equation~\eqref{eq:13.2.10} reduces to our previous criterion
\(\pi_1(B)=\pi_1(\Lambda)=Z_2\). \emph{If \(\pi_1(B)\) cannot be expressed in the
form~\eqref{eq:13.2.10}, the notion of spinors cannot be defined on \(M\).} On the
other hand, if \(\pi_1(B)\) is of the form~\eqref{eq:13.2.10}, we may construct a
(nonuniversal) covering space, \(\widehat B\), of \(B\) by defining two closed
curves through a point \(b\in B\) to be equivalent if their composition is homotopic
to a closed curve of the form \((e_1,g_2)\). Then \(\widehat B\) has the natural
structure of a principal \(\mathrm{SL}(2,\mathbb C)\) bundle over \(M\), and one can
proceed to define spinors on \(M\) in the same manner as in the case where \(M\) is
simply connected. Again---assuming that \(M\) is orientable and time
orientable---spinors can be defined [i.e., \(\pi_1(B)\) is of the
form~\eqref{eq:13.2.10}] if and only if the second Stiefel--Whitney class of \(M\)
vanishes. Furthermore, if in addition \(M\) is noncompact, spinors can be defined if
and only if \(M\) is parallelizable (Geroch 1968a; see also Clarke 1971).

One further point is worthy of note in the case where \(M\) is not simply connected.
The decomposition of \(\pi_1(B)\) as a direct product of subgroups of the
form~\eqref{eq:13.2.10} need not be unique, and each distinct decomposition gives
rise to a distinct notion of spinors on \(M\). The reason is as follows. Each
element of \(\pi_1(B)\) corresponds to an equivalence class of transports of tetrads
around closed curves in \(M\). When we express \(\pi_1(B)\) in the
form~\eqref{eq:13.2.10}, we, in effect, state that tetrad transports of the form
\((e_1,g_2)\) correspond to no net rotation of the tetrad, while those of the form
\((a_1,g_2)\)---where \(a_1\) is the non-identity element of
\(\pi_1(\Lambda)=Z_2\)---correspond to a \(2\pi\)-rotation of the tetrad. Suppose,
now, that \(\pi_1(M)\) has a normal subgroup\footnote{A subgroup \(H\) of a group
\(G\) is said to be \emph{normal} if \(ghg^{-1}\in H\) for all \(g\in G\) and
\(h\in H\).} \(H\) with factor group\footnote{For a normal subgroup \(H\), the
orbits obtained by the natural left action of \(H\) on \(G\)---called \emph{right
cosets} of \(H\)---can be given a group structure via the composition law
\(c_1c_2=c\), where \(c\) is the coset of \(H\) containing \(g_1g_2\), where
\(g_1\in c_1\) and \(g_2\in c_2\). (It is necessary that \(H\) be a normal subgroup
in order that this composition law be independent of the choice of representative
elements \(g_1,g_2\).) This group of cosets is called the \emph{factor group} of
\(H\) and is denoted \(G/H\).} isomorphic to \(Z_2\). Then we can define \(G_2'\)
to be the subset of \(\pi_1(B)\) consisting of all elements either of the form
\((e,h)\) or of the form \((a,bh)\), where \(h\in H\) and \(b\notin H\). Then it is
easy to check that \(G_2'\) is a subgroup of \(\pi_1(B)\) isomorphic to \(\pi_1(M)\)
and that \(\pi_1(B)\) is of the form \(G_1\times G_2'\). We may use this
decomposition to define the covering space \(\widehat B'\) and use \(\widehat B'\)
to define spinors. With this new decomposition of \(\pi_1(B)\), closed curves in
\(B\) of the form \((a,bh)\) now correspond to no net rotation of the tetrad in
\(M\), whereas curves of the form \((e,bh)\) now correspond to a
\(2\pi\)-rotation. Spinors on \(M\) defined with respect to \(\widehat B\) will not
change sign if transported around a closed curve \(bh\in\pi_1(M)\) in \(M\) in such a
way that the components of its null flag~\eqref{eq:13.1.37} with respect to the
basis associated with \((e,bh)\in\pi_1(B)\) remain constant during the transport.
Spinors defined with respect to \(\widehat B'\) will change sign under the same
transport. Thus, depending on the group structure of \(\pi_1(M)\), there may be
several\footnote{The number of distinct ways of defining spinors equals the number
of normal subgroups of \(\pi_1(M)\) with factor group \(Z_2\). This, in turn, is
equal to the number of generators of the cohomology group \(H^1(M;Z_2)\); see Isham
(1978) for further discussion.} distinct ways of defining spinors in the sense that
there may exist freedom to decide whether or not a spinor changes sign when
transported around certain closed curves in \(M\). For example, if \(\pi_1(M)=Z_2\),
then if spinors can be defined at all, there exist two inequivalent definitions. On
the other hand, if \(\pi_1(M)=Z_3\) (i.e., the cyclic group consisting of three
elements), then the definition of spinors is unique.

We turn, now, to the definition of derivatives of spinors in curved spacetime. As
mentioned above, in curved spacetime there is no natural identification of the
spinor spaces at different points, just as the tangent spaces at different points
cannot be identified in a natural way. However, in chapter~3 we introduced the
notion of a derivative operator on ordinary tensor fields and found that there is a
unique derivative operator, \(\nabla_a\), associated with a metric \(g_{ab}\) via
the requirement \(\nabla_a g_{bc}=0\). Furthermore, given a derivative operator, we
obtain a notion of parallel transport, i.e., a curve-dependent identification of the
tangent spaces at different points. This notion of parallel transport of tensor
fields obtained from \(\nabla_a\) gives rise to a unique notion of parallel transport
of spinors,\footnote{This can be demonstrated more systematically by reformulating
the notion of parallel transport in terms of the theory of connections on a principal
fiber bundle and its associated bundles (see, e.g., Bishop and Crittenden 1964). The
connection on \(B\) gives rise to a connection on \(\widehat B\), which in turn
yields a connection on \(S(M)\).} namely, a spinor will be said to be parallelly
transported along a curve if its null flag~\eqref{eq:13.1.37} is parallelly
transported and, in addition, the spinor itself does not change sign
discontinuously. (Recall that the continuity of a spinor field is a well-defined
notion.) This notion of parallel transport then can be used to take derivatives of
spinors along a curve \(\gamma\)---since a spinor at \(\gamma(t)\) now can be
compared with a spinor at \(\gamma(t+\delta t)\) via parallel transport---and this,
in turn, uniquely gives rise to a notion of a spinor derivative operator
\(\nabla_{AA'}\) taking spinorial tensor fields of type \((k,l;k',l')\) into
spinorial tensor fields of type \((k,l+1;k',l'+1)\). This derivative operator
\(\nabla_{AA'}\) will satisfy the analogs of properties (1)--(3) listed in
section~3.1 for a derivative operator on ordinary tensor fields where, in property
(3), contraction now is with respect to spinor indices. Furthermore, when applied
to ordinary tensor fields, \(\nabla_{AA'}\) will agree with the derivative operator
\(\nabla_a\) associated with \(g_{ab}\), which implies, in particular, that
properties (4) and (5) are satisfied by \(\nabla_{AA'}\). In addition,
\(\nabla_{AA'}\) satisfies two further conditions: (i) for all spinorial tensor
fields \(\psi\) we have \(\overline{\nabla\psi}=\nabla\overline\psi\), i.e.,
\(\nabla_{AA'}\) is \lq\lq real,\rq\rq\ and (ii)
\(\nabla_{AA'}\epsilon_{BC}=0=\nabla_{AA'}\overline\epsilon_{B'C'}\). Note that
although the derivative operator \(\nabla_a\) satisfying \(\nabla_a g_{bc}=0\) is
but one of a wide class of derivative operators on tensor fields, most of these
other derivative operators cannot be generalized to apply to spinor fields, since
for these notions of derivative the parallel transport of a null flag will not, in
general, retain the form of a null flag.

If we take two spinor derivatives of a dual spinor field \(\alpha_C\) and
antisymmetrize over the derivative indices, the same argument used to define the
Riemann tensor in chapter~3 shows that the resulting spinorial tensor field at a
point \(x\in M\) depends only on the value of \(\alpha_C\) at \(x\). Thus, there
exists a spinorial tensor field \(\chi_{AA'BB'C}{}^D\) such that for all dual
spinor fields \(\alpha_C\), we have
\begin{equation}
  \bigl(\nabla_{AA'}\nabla_{BB'}-\nabla_{BB'}\nabla_{AA'}\bigr)\alpha_C
  =\chi_{AA'BB'C}{}^D\alpha_D.
  \tag{13.2.11}\label{eq:13.2.11}
\end{equation}
We apply this commutator to the covector
\(\alpha_C\overline\alpha_{C'}\), using the Leibniz rule and the reality property
(i) of \(\nabla_{AA'}\). Comparison with formula~(3.2.3) for all
\(\alpha_C\overline\alpha_{C'}\) establishes that
\begin{equation}
  R_{AA'BB'CC'}{}^{DD'}
  =\chi_{AA'BB'C}{}^D\overline\epsilon_{C'}{}^{D'}
   +\overline\chi_{AA'BB'C'}{}^{D'}\epsilon_C{}^D,
  \tag{13.2.12}\label{eq:13.2.12}
\end{equation}
where \(R_{AA'BB'CC'}{}^{DD'}\) is the spinorial equivalent of the Riemann tensor.
Arguments similar to those which led to equation~(3.2.12) show that, for an
arbitrary \lq\lq down index\rq\rq\ spinorial tensor field,
\(\alpha_{C_1\cdots C_lC_1'\cdots C_{l'}'}\), we have
\begin{align}
  &\bigl(\nabla_{AA'}\nabla_{BB'}-\nabla_{BB'}\nabla_{AA'}\bigr)
    \alpha_{C_1\cdots C_lC_1'\cdots C_{l'}'}\notag\\
  &\quad={}\sum_{i=1}^l\chi_{AA'BB'C_i}{}^D
    \alpha_{C_1\cdots D\cdots C_lC_1'\cdots C_{l'}'}\notag\\
  &\qquad{}+\sum_{i=1}^{l'}\overline\chi_{AA'BB'C_i'}{}^{D'}
    \alpha_{C_1\cdots C_lC_1'\cdots D'\cdots C_{l'}'}.
  \tag{13.2.13}\label{eq:13.2.13}
\end{align}
The generalization of equation~\eqref{eq:13.2.13} to spinorial tensor fields with
\lq\lq up indices\rq\rq\ is easily obtained by raising indices on both sides of
equation~\eqref{eq:13.2.13} with \(\epsilon^{EC}\) and
\(\overline\epsilon^{E'C'}\).

The antisymmetry of \(R_{abcd}\) in its last two indices implies via
equations~\eqref{eq:13.2.12} and~\eqref{eq:13.1.46} that
\(\chi_{AA'BB'CD}\) is symmetric in \(C\) and \(D\). Since \(R_{abcd}\) is
antisymmetric in its first two indices, the argument which led to
equation~\eqref{eq:13.1.50} implies further that \(\chi_{AA'BB'CD}\) can be
expressed in the form
\begin{equation}
  \chi_{AA'BB'CD}
  =\Lambda_{ABCD}\overline\epsilon_{A'B'}
   +\Phi_{A'B'CD}\epsilon_{AB},
  \tag{13.2.14}\label{eq:13.2.14}
\end{equation}
where
\begin{equation}
  \Lambda_{ABCD}=\Lambda_{(AB)(CD)}
  \tag{13.2.15}\label{eq:13.2.15}
\end{equation}
and
\begin{equation}
  \Phi_{A'B'CD}=\Phi_{(A'B')(CD)}.
  \tag{13.2.16}\label{eq:13.2.16}
\end{equation}
Substituting these results in equation~\eqref{eq:13.2.12}, we find that the
Riemann tensor symmetry \(R_{abcd}=R_{cdab}\) implies that \(\Phi_{A'B'CD}\) is
real:
\begin{equation}
  \overline\Phi_{ABC'D'}=\Phi_{C'D'AB},
  \tag{13.2.17}\label{eq:13.2.17}
\end{equation}
and \(\Lambda_{ABCD}\) satisfies
\begin{equation}
  \Lambda_{ABCD}=\Lambda_{CDAB}.
  \tag{13.2.18}\label{eq:13.2.18}
\end{equation}
Because of equation~\eqref{eq:13.2.18},
\(\epsilon^{AC}\Lambda_{ABCD}\) is antisymmetric in \(B\) and \(D\) and hence must
be a multiple of \(\epsilon_{BD}\). We define \(\Psi_{ABCD}\) by
\begin{equation}
  \Psi_{ABCD}
  =\Lambda_{ABCD}-\Lambda(\epsilon_{AC}\epsilon_{BD}+\epsilon_{BC}\epsilon_{AD}),
  \tag{13.2.19}\label{eq:13.2.19}
\end{equation}
where
\begin{equation}
  \Lambda=\frac16\epsilon^{AC}\epsilon^{BD}\Lambda_{ABCD}.
  \tag{13.2.20}\label{eq:13.2.20}
\end{equation}
Then by construction \(\Psi_{ABCD}\) satisfies
\begin{equation}
  \epsilon^{AC}\Psi_{ABCD}=0.
  \tag{13.2.21}\label{eq:13.2.21}
\end{equation}
Together with the other symmetries of \(\Lambda_{ABCD}\), this implies that
\(\Psi_{ABCD}\) is totally symmetric,
\begin{equation}
  \Psi_{ABCD}=\Psi_{(ABCD)}.
  \tag{13.2.22}\label{eq:13.2.22}
\end{equation}
Finally, a calculation using equation~\eqref{eq:13.1.36} establishes that the
Riemann tensor symmetry \(R_{a[bcd]}=0\) implies the reality of \(\Lambda\),
\begin{equation}
  \overline\Lambda=\Lambda.
  \tag{13.2.23}\label{eq:13.2.23}
\end{equation}

Thus, we have found that
\begin{equation}
  \chi_{AA'BB'CD}
  =\Psi_{ABCD}\overline\epsilon_{A'B'}
   +\Phi_{A'B'CD}\epsilon_{AB}
   +\Lambda(\epsilon_{AC}\epsilon_{BD}+\epsilon_{BC}\epsilon_{AD})
    \overline\epsilon_{A'B'},
  \tag{13.2.24}\label{eq:13.2.24}
\end{equation}
and hence we have obtained the following spinorial decomposition of the Riemann
tensor:
\begin{align}
  R_{AA'BB'CC'DD'}
  ={}&\Psi_{ABCD}\overline\epsilon_{A'B'}\overline\epsilon_{C'D'}
    +\Phi_{A'B'CD}\epsilon_{AB}\overline\epsilon_{C'D'}\notag\\
   &{}+\Lambda(\epsilon_{AC}\epsilon_{BD}+\epsilon_{BC}\epsilon_{AD})
      \overline\epsilon_{A'B'}\overline\epsilon_{C'D'}
    +\text{C.C.},
  \tag{13.2.25}\label{eq:13.2.25}
\end{align}
where C.C.\ denotes the complex conjugate of the preceding terms. To obtain the
Ricci tensor, we contract over \(B\) and \(D\) and over \(B'\) and \(D'\). We obtain
\begin{equation}
  R_{AA'CC'}=-2\Phi_{A'C'AC}
    +6\Lambda\overline\epsilon_{A'C'}\epsilon_{AC}.
  \tag{13.2.26}\label{eq:13.2.26}
\end{equation}
This shows that \(-2\Phi_{A'C'AC}\) is just the spinor equivalent of the trace-free
Ricci tensor \(R_{ab}+\frac14Rg_{ab}\), where now we define
\(R=-R_{ab}g^{ab}\) in order to compensate for our change in metric signature and
thereby agree with our previous definition given in chapter~3. (The definitions of
\(R_{abcd}\) and \(R_{ab}\) are unaffected by the change of metric signature made
in this chapter.) Furthermore, equation~\eqref{eq:13.2.26} shows that
\(\Lambda=-R/24\). Using these results and comparing the decomposition
\eqref{eq:13.2.25} with that of~(3.2.28), we find that the spinor equivalent,
\(C_{AA'BB'CC'DD'}\), of the Weyl tensor \(C_{abcd}\) is simply
\begin{equation}
  C_{AA'BB'CC'DD'}
  =\Psi_{ABCD}\overline\epsilon_{A'B'}\overline\epsilon_{C'D'}
   +\overline\Psi_{A'B'C'D'}\epsilon_{AB}\epsilon_{CD}.
  \tag{13.2.27}\label{eq:13.2.27}
\end{equation}
We refer to \(\Psi_{ABCD}\) as the \emph{Weyl spinor}.

The decomposition~\eqref{eq:13.2.24} of \(\chi_{AA'BB'C}{}^D\) allows us to
reexpress equation~\eqref{eq:13.2.11} in such a way as to separate the effect of
Ricci and Weyl curvature on spinors. If we contract equation~\eqref{eq:13.2.11}
with \(\epsilon^{AB}\), we find
\begin{equation}
  \nabla_{(A'}{}^A\nabla_{B')A}\alpha_C
  =\Phi_{A'B'C}{}^D\alpha_D,
  \tag{13.2.28}\label{eq:13.2.28}
\end{equation}
where we used the identity~\eqref{eq:13.1.46} to obtain the left-hand side of this
equation and formula~\eqref{eq:13.2.24} to obtain the right-hand side. On the other
hand, if we contract equation~\eqref{eq:13.2.11} with
\(\overline\epsilon^{A'B'}\), we get
\begin{equation}
  \nabla_{A'(A}\nabla_{B)}{}^{A'}\alpha_C
  =\Lambda_{ABC}{}^D\alpha_D
  =\Psi_{ABC}{}^D\alpha_D-2\Lambda\epsilon_{C(A}\alpha_{B)}.
  \tag{13.2.29}\label{eq:13.2.29}
\end{equation}
The two equations~\eqref{eq:13.2.28} and~\eqref{eq:13.2.29} are equivalent to
equation~\eqref{eq:13.2.11}. For future reference, we note that the generalization
of equation~\eqref{eq:13.2.29} to an \(n\) index spinor
\(\phi_{A_1\cdots A_n}\) is
\begin{align}
  \nabla_{A'(A}\nabla_{B)}{}^{A'}\phi_{C_1\cdots C_n}
  ={}&\sum_{i=1}^n\bigl\{
    \Psi_{ABC_i}{}^D\phi_{C_1\cdots D\cdots C_n}\notag\\
   &\qquad{}-2\Lambda\epsilon_{C_i(A}
     \phi_{C_1\cdots C_{i-1}B)C_{i+1}\cdots C_n}
    \bigr\}.
  \tag{13.2.30}\label{eq:13.2.30}
\end{align}
Note also that the general identity~\eqref{eq:13.1.52} implies that
\begin{equation}
  \nabla_{A'[A}\nabla_{B]}{}^{A'}=\frac12\epsilon_{AB}\Box,
  \tag{13.2.31}\label{eq:13.2.31}
\end{equation}
where \(\Box=\nabla_{AA'}\nabla^{AA'}\).

We now are in a position to explain the spinor motivation of the Newman--Penrose
(1962) formalism mentioned in section~3.4. Instead of choosing an orthonormal basis
of the tangent space at each point, we choose a basis
\(\xi_0^A=o^A,\ \xi_1^A=\iota^A\) of spinor space at each point, normalized so that
\begin{equation}
  o_A\iota^A=1.
  \tag{13.2.32}\label{eq:13.2.32}
\end{equation}
Instead of defining the 24 real Ricci rotation coefficients by equation~(3.4.14), we
define the 12 complex \emph{spin coefficients} by
\begin{equation}
  \gamma_{\Gamma\Delta'\Sigma\Lambda}
  =(\xi_\Gamma)^A(\overline\xi_{\Delta'})^{A'}(\xi_\Sigma)_B
    \nabla_{AA'}(\xi_\Lambda)^B.
  \tag{13.2.33}\label{eq:13.2.33}
\end{equation}
Here \(\Gamma,\Delta',\Sigma,\Lambda\) take the values 0, 1 and using the Leibniz
rule and equation~\eqref{eq:13.2.32} one verifies that
\(\gamma_{\Gamma\Delta'\Sigma\Lambda}\) is symmetric in \(\Sigma\) and \(\Lambda\),
so equation~\eqref{eq:13.2.33} indeed defines 12 independent complex quantities.
For the specific notation customarily used for each spin coefficient, see Newman and
Penrose (1962). Finally, instead of using equation~(3.4.17) to express the tetrad
basis components of the Riemann tensor in terms of the Ricci rotation coefficients,
we use equations~\eqref{eq:13.2.11} and~\eqref{eq:13.2.24} to express the spinor
basis components of \(\Psi_{ABCD}\), \(\Phi_{A'B'CD}\), and \(\Lambda\) in terms of
the spin coefficients. In this way, we obtain equations equivalent to
equation~(3.4.21) of chapter~3, but the explicit form of these equations written out
term by term (see eq.~[4.2] of Newman and Penrose 1962) is far simpler in appearance
than equation~(3.4.21) would be if a separate letter were used for each Ricci
rotation coefficient and each curvature component.

How can one evaluate the spin coefficients defined by equation~\eqref{eq:13.2.33}?
More generally, how does one compute derivatives, \(\nabla_{AA'}\), of arbitrary
spinorial tensors? Recall that the spinor derivative operator \(\nabla_{AA'}\)
corresponding to the tensor derivative operator \(\nabla_a\) satisfying
\(\nabla_a g_{bc}=0\) is the only sensible derivative operator for spinors, so no
analog of an \lq\lq ordinary derivative\rq\rq\ operator \(\partial_{AA'}\) exists.
Thus, there is no analog for spinors of the coordinate basis methods used for
tensors. However, we can
evaluate the spin coefficients~\eqref{eq:13.2.33} by using the formula (Friedman,
unpublished)
\begin{equation}
  \gamma_{AA'\Sigma\Lambda}
  \equiv(\xi_\Sigma)_B\nabla_{AA'}(\xi_\Lambda)^B
  =\frac12\sum_{\Gamma',\Delta'=0}^1
    \overline\epsilon^{\Gamma'\Delta'}(\overline\xi_{\Gamma'})_{B'}
    (\xi_\Sigma)_B\nabla_{AA'}
    \bigl[(\xi_\Lambda)^B(\overline\xi_{\Delta'})^{B'}\bigr].
  % SOURCE ERRATUM: printed tag (10.2.34); corrected to (13.2.34).
  \tag{13.2.34}\label{eq:13.2.34}
\end{equation}
The term in square brackets is a vector quantity, i.e.,
\begin{equation}
  l^{AA'}=\iota^A\overline\iota^{A'},
  \tag{13.2.35}\label{eq:13.2.35}
\end{equation}
and
\begin{equation}
  n^{AA'}=o^A\overline o^{A'}
  \tag{13.2.36}\label{eq:13.2.36}
\end{equation}
are null vectors, while
\begin{equation}
  m^{AA'}=\iota^A\overline o^{A'}
  \tag{13.2.37}\label{eq:13.2.37}
\end{equation}
and
\begin{equation}
  \overline m^{AA'}=o^A\overline\iota^{A'}
  \tag{13.2.38}\label{eq:13.2.38}
\end{equation}
are (complex) linear combinations of the spacelike unit vectors
\begin{equation}
  x^{AA'}=\frac{1}{\sqrt2}
    \bigl(\overline\iota^{A'}o^A+\iota^A\overline o^{A'}\bigr),
  \tag{13.2.39}\label{eq:13.2.39}
\end{equation}
\begin{equation}
  y^{AA'}=\frac{\ii}{\sqrt2}
    \bigl(\overline\iota^{A'}o^A-\iota^A\overline o^{A'}\bigr).
  \tag{13.2.40}\label{eq:13.2.40}
\end{equation}
Thus, the spin coefficients can be evaluated in terms of
\(\nabla_a l^b\), \(\nabla_a n^b\), \(\nabla_a x^b\), and \(\nabla_a y^b\), which
can be calculated by standard methods. Indeed, the entire spinor calculus can be
reexpressed as a tetrad calculus using the complex null tetrad
\(l^a,n^a,m^a,\overline m^a\). Once the spin coefficients have been found, the
derivative of an arbitrary spinorial tensor field \(\psi\) can be evaluated by first
expanding \(\psi\) in a basis formed out of \(\iota^A,o^A\) and their complex
conjugates. Then, using the Leibniz rule, we can express \(\nabla\psi\) in terms of
derivatives of the (scalar) basis components and the spin coefficients.

We illustrate, now, the utility of spinor methods by deriving the algebraic
classification of the Weyl tensor (Penrose 1960). In section~7.3 we asserted that in
algebraically general spacetimes the Weyl tensor admits four distinct principal null
directions satisfying equation~(7.3.1), whereas in algebraically special spacetimes
some of these null directions coincide and satisfy the stronger conditions listed in
Table~7.1. We did not attempt to prove these assertions in chapter~7 because a
direct proof by tensor methods is relatively difficult. However, the algebraic
classification of the Weyl tensor by spinor methods is remarkably simple. Consider
the Weyl spinor \(\Psi_{ABCD}\) and fix a basis \(\iota^A,o^A\) of spinor space with
\(o_A\iota^A=1\) and with \(\iota^A\) chosen so that
\begin{equation}
  \Psi_{ABCD}\iota^A\iota^B\iota^C\iota^D=1.
  \tag{13.2.41}\label{eq:13.2.41}
\end{equation}
Let \(z\in\mathbb C\), set
\begin{equation}
  \alpha^A=z\iota^A+o^A,
  \tag{13.2.42}\label{eq:13.2.42}
\end{equation}
and consider the quantity
\begin{equation}
  f(z)=\Psi_{ABCD}\alpha^A\alpha^B\alpha^C\alpha^D.
  \tag{13.2.43}\label{eq:13.2.43}
\end{equation}
Since \(f\) is a fourth degree polynomial in \(z\) and by
equation~\eqref{eq:13.2.41} the coefficient of \(z^4\) is equal to 1, we can factor
it into the form
\begin{equation}
  f(z)=(z-c_1)(z-c_2)(z-c_3)(z-c_4).
  \tag{13.2.44}\label{eq:13.2.44}
\end{equation}
Now, for \(i=1,2,3,4\), let
\begin{equation}
  (\kappa_i)^A=o^A+c_i\iota^A.
  \tag{13.2.45}\label{eq:13.2.45}
\end{equation}
Then we have
\begin{equation}
  z-c_i=(\kappa_i)_A\alpha^A,
  \tag{13.2.46}\label{eq:13.2.46}
\end{equation}
and hence we have established that there exist four spinors \((\kappa_i)^A\), called
\emph{principal spinors}, such that the equation
\begin{equation}
  \Psi_{ABCD}\alpha^A\alpha^B\alpha^C\alpha^D
  =(\kappa_1)_A(\kappa_2)_B(\kappa_3)_C(\kappa_4)_D
    \alpha^A\alpha^B\alpha^C\alpha^D
  \tag{13.2.47}\label{eq:13.2.47}
\end{equation}
holds for all \(\alpha^A\) of the form~\eqref{eq:13.2.42}, and, consequently, for
all spinors \(\alpha^A\). However, since the Weyl spinor is totally symmetric (see
eq.~[13.2.22]), this can be true if and only if
\begin{equation}
  \Psi_{ABCD}
  =(\kappa_1)_{(A}(\kappa_2)_B(\kappa_3)_C(\kappa_4)_{D)}.
  \tag{13.2.48}\label{eq:13.2.48}
\end{equation}
Thus, we obtain a general decomposition of the Weyl spinor as the symmetrized
product of four principal spinors. Note that \(\kappa^A\) is a principal spinor if
and only if
\begin{equation}
  \Psi_{ABCD}\kappa^A\kappa^B\kappa^C\kappa^D=0.
  \tag{13.2.49}\label{eq:13.2.49}
\end{equation}
Furthermore, \(\kappa^A\) is 2, 3, or 4 times repeated---i.e., it appears 2, 3, or
4 times in the decomposition~\eqref{eq:13.2.48}---if and only if, respectively,
\begin{align}
  \Psi_{ABCD}\kappa^A\kappa^B\kappa^C&=0
  &&\text{(two times)}, \tag{13.2.50}\label{eq:13.2.50}\\
  \Psi_{ABCD}\kappa^A\kappa^B&=0
  &&\text{(three times)}, \tag{13.2.51}\label{eq:13.2.51}\\
  \Psi_{ABCD}\kappa^A&=0
  &&\text{(four times)}. \tag{13.2.52}\label{eq:13.2.52}
\end{align}
If \(\kappa^A\) is a principal spinor, the null vector
\begin{equation}
  k^{AA'}=\kappa^A\overline\kappa^{A'}
  \tag{13.2.53}\label{eq:13.2.53}
\end{equation}
is called a principal null vector. The conditions~\eqref{eq:13.2.49}--\eqref{eq:13.2.52}
on \(\kappa^A\) now can be translated to the conditions on \(k^a\) listed in
Table~7.1 (problem~5). Thus, we obtain the algebraic classification of
section~7.3.

We conclude this section by examining the generalization to curved spacetime of the
equations in Minkowski spacetime for fields of mass \(m\) and spin \(s\). The most
natural generalization of these equations is obtained by the \lq\lq minimal
coupling\rq\rq\ prescription of replacing \(\partial_{AA'}\) everywhere by
\(\nabla_{AA'}\). However, for \(m>0\), this prescription gives inequivalent
results when applied to equation~\eqref{eq:13.1.59} as opposed to
equations~\eqref{eq:13.1.60} and~\eqref{eq:13.1.61}. The curved spacetime version
of equation~\eqref{eq:13.1.59} has a well posed initial value formulation for all
\(s\), but the current~\eqref{eq:13.1.62} is no longer conserved. On the other
hand, for \(s=\frac12\) the equations
\begin{equation}
  \nabla_{AA'}\phi^A=\frac{m}{\sqrt2}\sigma_{A'},
  \tag{13.2.54}\label{eq:13.2.54}
\end{equation}
\begin{equation}
  \nabla^{AA'}\sigma_{A'}=-\frac{m}{\sqrt2}\phi^A
  \tag{13.2.55}\label{eq:13.2.55}
\end{equation}
have a well posed initial value formulation (see problem~8), and the
current~\eqref{eq:13.1.62} generalized to curved spacetime is conserved. Hence, we
adopt equations~\eqref{eq:13.2.54} and~\eqref{eq:13.2.55} as the generalization to
curved spacetime for the Dirac equation. Similar results hold for \(s=1\). However,
when \(s>1\), the curved spacetime versions of equations~\eqref{eq:13.1.60} and
\eqref{eq:13.1.61} do not have a well posed initial value formulation (Buchdahl
1962). Similarly, the curved spacetime spin \(s\) equation for \(m=0\),
\begin{equation}
  \nabla_{A_1A_1'}\phi^{A_1\cdots A_n}=0,
  \tag{13.2.56}\label{eq:13.2.56}
\end{equation}
has a well posed initial value formulation for \(s=\frac12\) and for \(s=1\) (where
it is equivalent to Maxwell's equations in curved spacetime), but, as we shall see
below, it fails to have a well posed initial value formulation for \(s>1\). Thus, the
natural curved spacetime generalization of the Minkowski equations for \(s>1\) do
not yield physically viable models for fields in curved spacetime.

To prove that equation~\eqref{eq:13.2.56} does not have a well posed initial value
formulation when \(s>1\), we note first that equations~\eqref{eq:13.2.56}
and~\eqref{eq:13.2.31} imply that \(\phi^{A_1\cdots A_n}\) must satisfy the wave
equation,
\begin{equation}
  \Box\phi^{A_1\cdots A_n}=0.
  \tag{13.2.57}\label{eq:13.2.57}
\end{equation}
However, unlike the situation in flat spacetime, if equation~\eqref{eq:13.2.57} is
satisfied everywhere in spacetime and equation~\eqref{eq:13.2.56} holds on an
initial surface, this does not imply that equation~\eqref{eq:13.2.56} holds
everywhere. Indeed, contracting equation~\eqref{eq:13.2.56} with
\(\nabla_{A_1'A_2}\) and using the fact that
\(\phi^{A_1\cdots A_n}=\phi^{(A_1\cdots A_n)}\), we obtain
\begin{align}
  0
  &={}\nabla_{A_1'A_2}\nabla_{A_1}{}^{A_1'}
    \phi^{A_1A_2\cdots A_n}\notag\\
  &={}\nabla_{A_1'(A_2}\nabla_{A_1)}{}^{A_1'}
    \phi^{A_1A_2\cdots A_n}\notag\\
  &={}(n-2)\Psi_{A_1A_2C(A_3}
    \phi^{\lvert A_1A_2C\rvert A_4\cdots A_n)}.
  \tag{13.2.58}\label{eq:13.2.58}
\end{align}
where equation~\eqref{eq:13.2.30} was used in the last line. Equation~\eqref{eq:13.2.58}
is a purely algebraic condition which must be satisfied by
\(\phi^{A_1\cdots A_n}\) throughout spacetime. Note that it is \emph{not}
preserved under evolution by equation~\eqref{eq:13.2.57}; i.e., even if
equation~\eqref{eq:13.2.58} and its time derivative hold for data for
\(\phi^{A_1\cdots A_n}\) on an initial surface, there is no reason why
equation~\eqref{eq:13.2.58} will hold for the solution of equation~\eqref{eq:13.2.57}
evolved from these data. Thus, few, if any, solutions of equation~\eqref{eq:13.2.56}
will exist in a general curved spacetime. Thus, for \(s>1\) there is no natural
generalization to curved spacetime of the notion of a \lq\lq pure\rq\rq\ massless
spin \(s\) field.
